\documentclass[a4paper, 12pt]{article}
% \documentclass[a4paper, 12pt]{report}

\usepackage[T1]{fontenc}
\usepackage[utf8]{inputenc}
\usepackage[brazilian]{babel}
\usepackage{graphicx}
\usepackage{float}
\usepackage{hyperref}
\usepackage{listings}

\title{\textbf{Computação de Alto Desempenho}\\Trabalho 1\\\vspace{10cm}}
\author{
Pedro Henrique Grave Lima - 120019143
}
\date{\today}

\begin{document}
\maketitle
\newpage

\section*{Introdução}
Um único `relatorio.pdf` na raiz do arquivo compactado, contendo:

\begin{itemize}
    \item Uma capa curta (nome, DRE, data).
    \item Para cada exercício: uma descrição breve da sua abordagem, os comandos usados para compilar e executar, a saída, e qualquer discussão pedida pelo enunciado.
    \item Para o Ex. 1(a): a observação das execuções múltiplas (PIDs distintos, coluna de CPU mudando).
    \item Para o Ex. 2(a): a saída de verificação com $p = 3, n = 10$.
    \item Para o Ex. 2(b): a tabela de $T_p$, $S_p$, $E_p$ e a discussão de quando $E_p$ cai abaixo de 0.8.
    \item Para o Ex. 2(c) (se feito): o gráfico log–log de convergência com a inclinação medida.
    \item Para o Ex. 3: o erro relativo contra a soma serial e uma comparação de contagens de linhas com a versão da Aula 6 usando coletivas.
    \item Para o Ex. 5: a discussão de quando Allgather é preferível a Gather.
    \item Para o Ex. 6: a comparação entre a versão com tipo derivado e a versão com três Bcasts.
\end{itemize}

LaTeX puro, Markdown convertido para PDF, ou qualquer ferramenta que produza um PDF limpo é aceito. Gráficos feitos à mão \textbf{não} são aceitos; use matplotlib, gnuplot, ou equivalente.

Neste trabalho, os códigos serão executados no ambiente WSL2 em um computador equipado com um processador \href{https://www.intel.com.br/content/www/br/pt/products/sku/236789/intel-core-i7-processor-14700kf-33m-cache-up-to-5-60-ghz/specifications.html}{Intel Core i7 14700KF}, que conta com um total de 28 threads.

\section{Exercício 1: Variantes do Hello World}
Ambos os códigos deste exercício foram desenvolvidos com foco em observar os identificadores de processos (PIDs) e como o MPI se comporta.

\subsection{Introspecção em nível de SO}
Na sua forma mais simples, vemos como o MPI dispara um processo único para cada thread solicitada. Neste caso pediremos 8 threads, argumento passado ao MPI em \verb|MPI_Init(&argc, &argv);|. Com isso, o programa abre 8 processos diferentes, cada um com seu próprio PID e alocado a uma thread de acordo com a disponibilidade naquele instante. Como existem diversos processos executando a todo momento, o núcleo onde o cada PID será executado pode ser considerado aleatório.

\begin{lstlisting}[language=Bash, caption=Exemplo de uma execução do programa]
    $ mpiexec -n 8 ./hello_os
    Hello from rank 1/8 -- PID = 6571, CPU = 8
    Hello from rank 2/8 -- PID = 6572, CPU = 1
    Hello from rank 3/8 -- PID = 6573, CPU = 14
    Hello from rank 4/8 -- PID = 6574, CPU = 12
    Hello from rank 5/8 -- PID = 6575, CPU = 4
    Hello from rank 7/8 -- PID = 6577, CPU = 2
    Hello from rank 6/8 -- PID = 6576, CPU = 0
    Hello from rank 0/8 -- PID = 6570, CPU = 4
\end{lstlisting}

Após executar diversas vezes, é possível observar que o PID cresce linearmente, conforme novos processos são abertos identificadores são criados individualmente.

\subsection{Variante em anel}



\end{document}
